\Askhsh{18566}{4}
\begin{ylh}{1}
    \item 10.3. Εμβαδόν βασικών ευθύγραμμων σχημάτων.
\end{ylh}
Δίνεται ορθογώνιο τρίγωνο $\text{ΑΒΓ}$ με \(\hat{\text{A}}=90^\circ\).
Με πλευρά την υποτείνουσα $\text{ΒΓ}$ και έξω από το τρίγωνο, γράφουμε το τετράγωνο $\text{ΒΓΔΕ}$.
Προεκτείνουμε την πλευρά $\text{ΒΑ}$ προς το $\text{Α}$ και παίρνουμε σημείο $\text{Ζ}$ τέτοιο ώστε $\text{ΒΖ} = \text{ΒΓ}$.
Από τα σημεία $\text{Γ}$ και $\text{Ζ}$ φέρουμε παράλληλες προς τα τμήματα $\text{ΒΖ}$ και $\text{ΒΓ}$ αντίστοιχα, που τέμνονται στο σημείο $\text{Η}$.
\begin{alist}
\item Να δικαιολογήσετε γιατί το τετράπλευρο $\text{ΒΓΗΖ}$ είναι ρόμβος και να βρείτε τις περιμέτρους του ρόμβου και του τετραγώνου. \monades{10}
\item Δίνονται οι ισχυρισμοί:
\begin{rlist}
\item Ισχυρισμός 1: «Ο ρόμβος και το τετράγωνο αφού έχουν ίσες περιμέτρους, θα έχουν και ίσα εμβαδά».
\item Ισχυρισμός 2: «Ο ρόμβος έχει μικρότερο εμβαδό από το τετράγωνο και μάλιστα, όπως έχουν κατασκευαστεί τα δύο τετράπλευρα δεν γίνεται να είναι ποτέ ισεμβαδικά».
\end{rlist}
Εξετάστε ποιος από τους 2 παραπάνω ισχυρισμούς είναι σωστός και να δικαιολογήσετε πλήρως την απάντησή σας. \monades{15}
\end{alist}